\section*{Capítulo \ref{ch:b3:forms} --- Formas diferenciais e
o teorema de Stokes}

\begin{solution}{exo:b3:forms:1}
Expandindo e matando os fatores repetidos:
\[
\omega\wedge\eta = (x\,\dd y - z\,\dd x)\wedge(\dd x +
y\,\dd z)
= -x\,\dd x\wedge\dd y + xy\,\dd y\wedge\dd z +
yz\,\dd z\wedge\dd x
\]
(usando $\dd y\wedge\dd x = -\dd x\wedge\dd y$ e $-zy\,\dd
x\wedge\dd z = yz\,\dd z\wedge\dd x$). A seguir, $\dd\omega =
\dd x\wedge\dd y - \dd z\wedge\dd x = \dd x\wedge\dd y +
\dd x\wedge\dd z$ e $\dd\eta = \dd y\wedge\dd z$. Por fim,
\[
\dd(\omega\wedge\eta) = y\,\dd x\wedge\dd y\wedge\dd z +
z\,\dd y\wedge\dd z\wedge\dd x = (y + z)\,\dd x\wedge\dd
y\wedge\dd z
\]
(o primeiro termo de $\omega\wedge\eta$ contribui com $\dd(-x)
\wedge\dd x\wedge\dd y = 0$; permutações cíclicas de três
fatores são pares). Verificação de Leibniz: $\dd\omega\wedge\eta =
(\dd x\wedge\dd y + \dd x\wedge\dd z)\wedge(\dd x + y\,\dd
z) = y\,\dd x\wedge\dd y\wedge\dd z$, e
$(-1)^1\omega\wedge\dd\eta = -(x\,\dd y - z\,\dd
x)\wedge\dd y\wedge\dd z = z\,\dd x\wedge\dd y\wedge\dd z$;
a soma bate.
\end{solution}

\begin{solution}{exo:b3:forms:2}
$\dd f = \sum_i\partial_if\,\dd x_i$ tem os coeficientes de
$\nabla f$. Para a forma trabalho,
\[
\dd\omega_F = (\partial_yF_3 - \partial_zF_2)\,\dd y\wedge\dd
z + (\partial_zF_1 - \partial_xF_3)\,\dd z\wedge\dd x +
(\partial_xF_2 - \partial_yF_1)\,\dd x\wedge\dd y =
\sigma_{\operatorname{curl}F},
\]
a forma \omterm{ex:b3:ode:planeclassification}{fluxo} do rotacional (reúna os seis termos de
$\sum_i\dd F_i\wedge\dd x_i$). Para a forma \omterm{ex:b3:ode:planeclassification}{fluxo},
$\dd\sigma_F = (\partial_xF_1 + \partial_yF_2 +
\partial_zF_3)\,\dd x\wedge\dd y\wedge\dd z$: a
divergência. Então $\dd^2f = 0$ se lê
$\sigma_{\operatorname{curl}\nabla f} = 0$, isto é,
$\operatorname{curl}\nabla f = 0$, e $\dd^2\omega_F = 0$
se lê $(\operatorname{div}\operatorname{curl}F)\,\dd
x\wedge\dd y\wedge\dd z = 0$: as duas identidades vetoriais são
uma única identidade, $\dd^2 = 0$, lida em dois graus.
\end{solution}

\begin{solution}{exo:b3:forms:3}
(a) Ser \omterm{def:b3:forms:closedexact}{fechada} é a simetria das derivadas parciais cruzadas:
$\partial_y(2xy + z^2) = 2x = \partial_x(x^2)$,
$\partial_z(2xy + z^2) = 2z = \partial_x(2xz)$,
$\partial_z(x^2) = 0 = \partial_y(2xz)$. O domínio $\R^3$
é estrelado: \omterm{def:b3:forms:closedexact}{exata} (\cref{thm:b3:forms:poincare}), com
primitiva $f = x^2y + xz^2$ (verifique $\dd f$).
(b) $\dfrac{x\,\dd x + y\,\dd y}{x^2 + y^2} =
\tfrac12\,\dd\log(x^2 + y^2)$: \omterm{def:b3:forms:closedexact}{exata} em todo
$\R^2\setminus\{0\}$ (logo \omterm{def:b3:forms:closedexact}{fechada}) --- a prima radial
da forma angular é inofensiva.
(c) Em $\{x > 0\}$ a forma é $\omega_\theta$, \omterm{def:b3:forms:closedexact}{fechada}
(\cref{exo:b3:forms:4}); o semiplano é convexo, de modo que ela é
\omterm{def:b3:forms:closedexact}{exata} ali e, de fato, $f = \arctan(y/x)$ satisfaz $\dd f
= \frac{-y\,\dd x + x\,\dd y}{x^2 + y^2}$. \omterm{def:b3:forms:closedexact}{Exata} no
semiplano, não \omterm{def:b3:forms:closedexact}{exata} no plano perfurado: o
obstáculo mora no buraco, não na fórmula.
\end{solution}

\begin{solution}{exo:b3:forms:4}
Ser \omterm{def:b3:forms:closedexact}{fechada}: com $\rho^2 = x^2 + y^2$,
\[
\partial_x\Bigl(\frac{x}{\rho^2}\Bigr) = \frac{\rho^2 -
2x^2}{\rho^4} = \frac{y^2 - x^2}{\rho^4} =
\partial_y\Bigl(\frac{-y}{\rho^2}\Bigr),
\]
de modo que $\dd\omega_\theta = \bigl(\partial_x(x/\rho^2) -
\partial_y(-y/\rho^2)\bigr)\,\dd x\wedge\dd y = 0$. Em
$\gamma(t) = (r\cos2\pi t, r\sin2\pi t)$:
\[
\gamma^*\omega_\theta = \frac{2\pi r^2(\cos^22\pi t +
\sin^22\pi t)}{r^2}\,\dd t = 2\pi\,\dd t,
\qquad\text{logo}\qquad
\int_\gamma\omega_\theta = 2\pi .
\] Se $\omega_\theta$ fosse
\omterm{def:b3:forms:closedexact}{exata}, essa integral se anularia
(\cref{prop:b3:forms:exactloop}): ela não é \omterm{def:b3:forms:closedexact}{exata}. Se
$\R^2\setminus\{0\}$ fosse estrelado em relação a algum $p$,
o lema de Poincaré (transladado a $p$) tornaria \omterm{def:b3:forms:closedexact}{exata} toda \omterm{def:b3:forms:closedexact}{forma
fechada} --- contradição. Diretamente: para qualquer $p
\neq 0$, o segmento de $p$ ao ponto $-p$ do
domínio passa por $0$: a propriedade de ser estrelado falha em todo
ponto.
\end{solution}

\begin{solution}{exo:b3:forms:5}
Para vetores $v_1, \dots, v_{n-1}$, expandir o determinante
pela primeira coluna dá
\[
\det(\nu, v_1, \dots, v_{n-1}) =
\sum_{i=1}^n(-1)^{i-1}\nu_i\,
\det\bigl(\text{linhas} \neq i \text{ de } (v_1, \dots,
v_{n-1})\bigr)
= \sigma_\nu(v_1, \dots, v_{n-1}),
\]
já que $(\dd x_1\wedge\dots\wedge\widehat{\dd
x_i}\wedge\dots\wedge\dd x_n)(v_1, \dots, v_{n-1})$ é
exatamente o $i$-ésimo menor obtido por supressão
(\cref{def:b3:forms:wedge}). Aplicando isso a $v_j =
\partial_j\gamma$: $(\gamma^*\sigma_\nu)(e_1, \dots,
e_{n-1}) = \det A$ com $A = (\nu\ \partial_1\gamma\ \cdots\
\partial_{n-1}\gamma)$. Ora, ${}^t\!AA$ é diagonal por blocos:
$\langle\nu, \nu\rangle = 1$ e $\langle\nu,
\partial_j\gamma\rangle = 0$ ($\nu$ normal, os
$\partial_j\gamma$ tangentes), de modo que $(\det A)^2 =
\det({}^t\!AA) = \det G$; e $\det A > 0$ para uma parametrização
direta (é isso que significa a \omterm{def:b3:forms:orientation}{orientação} por $\nu$):
$\gamma^*\sigma_\nu = \sqrt{\det G}\,\dd
u_1\wedge\dots\wedge\dd u_{n-1}$, o elemento de área de Gram. Para
$S^2$, $\nu(x) = x$ e a parametrização esférica
$\gamma(\theta, \varphi) = (\sin\varphi\cos\theta,
\sin\varphi\sin\theta, \cos\varphi)$ em
$(0,2\pi)\times(0,\pi)$ (direta; ela deixa de fora um meridiano, conjunto
que não carrega área): $\det G = \sin^2\varphi$, de modo que
$\int_{S^2}\sigma_\nu =
\int_0^{2\pi}\!\!\int_0^\pi\sin\varphi\,\dd\varphi\,
\dd\theta = 4\pi$.
\end{solution}

\begin{solution}{exo:b3:forms:6}
A $\omega$ dada é $\sigma_\nu$ para $\nu(x) = x$ em
$S^2$, de modo que (a) é o cálculo que acabamos de fazer:
$\int_{S^2}\omega = 4\pi$. (b) $\dd\omega = 3\,\dd
x\wedge\dd y\wedge\dd z$, e Stokes na bola unitária dá
$\int_{S^2}\omega = 3\operatorname{vol}(B^3)$: logo
$\operatorname{vol}(B^3) = \frac{4\pi}3$. Em geral, a
forma $\sigma = \sum_i(-1)^{i-1}x_i\,\dd
x_1\wedge\dots\wedge\widehat{\dd x_i}\wedge\dots\wedge\dd
x_n$ se restringe em $S^{n-1}$ à forma de área ($\nu = x$ no~\cref{exo:b3:forms:5}), $\dd\sigma = n\,\dd
x_1\wedge\dots\wedge\dd x_n$, e Stokes fornece
$\operatorname{area}(S^{n-1}) = n\operatorname{vol}(B^n)$
--- coerente com as fórmulas com função gama do~\cref{thm:b3:product:ballvolume}.
\end{solution}

\begin{solution}{exo:b3:forms:7}
(a) Stokes aplicado à forma \omterm{ex:b3:ode:planeclassification}{fluxo} $\sigma_F$ no
domínio \omterm{def:b3:topology:compact}{compacto} $\Omega$:
$\int_\Omega\operatorname{div}F\,\dd x =
\int_{\partial\Omega}\sigma_F$ (\cref{exo:b3:forms:2} para
o lado \omterm{def:b3:topology:interior}{interior}). Identifique o integrando de \omterm{def:b3:forms:boundary}{bordo}: para
$v_1, v_2$ tangentes num ponto de $\partial\Omega$, a
expansão pela primeira coluna do~\cref{exo:b3:forms:5} dá
$\sigma_F(v_1, v_2) = \det(F, v_1, v_2)$; escrevendo $F =
\langle F, \nu\rangle\nu + T$ com $T$ tangente, a coluna
$T$ é uma combinação do plano de $v_1, v_2$, de modo que $\det(T,
v_1, v_2) = 0$ e
$\sigma_F\vert_{\partial\Omega} = \langle F,
\nu\rangle\,\sigma_\nu = \langle F, \nu\rangle\,\dd S$: o
teorema da divergência, com $\nu$ a normal exterior
(normal-exterior-primeiro é exatamente a \omterm{def:b3:forms:orientation}{orientação} induzida).
(b) Stokes aplicado a $\omega_F$ na
superfície com \omterm{def:b3:forms:boundary}{bordo} $S$: $\dd\omega_F =
\sigma_{\operatorname{curl}F}$ se restringe a
$\langle\operatorname{curl}F, \nu\rangle\,\dd S$ pela mesma
identificação, ao passo que, na curva de \omterm{def:b3:forms:boundary}{bordo},
$\gamma^*\omega_F = \langle F\circ\gamma,
\gamma'\rangle\,\dd t$, isto é, $\oint\langle F,
\tau\rangle\,\dd\ell$. Semiesfera superior com $\nu$ exterior
(radial): no ponto do equador $p = (1,0,0)$, o
vetor exterior dentro da superfície é $-e_3$; completá-lo
em referenciais positivos mostra que o equador é percorrido no sentido
anti-horário visto de cima ($+e_3$): a regra da mão
direita, mesma convenção dos dois lados da identidade.
\end{solution}

\begin{solution}{exo:b3:forms:8}
Aplique o teorema da divergência (\cref{exo:b3:forms:7}(a),
cuja demonstração independe da dimensão) a $F = u\nabla v$:
$\operatorname{div}(u\nabla v) = u\,\Delta v + \langle\nabla
u, \nabla v\rangle$ e $\langle F, \nu\rangle =
u\,\partial_\nu v$: a primeira identidade. Trocar $u, v$ e
subtrair cancela o termo simétrico: a segunda. Se
$\Delta u = 0$ em $\Omega$ e $u = 0$ em $\partial\Omega$:
a primeira identidade com $v = u$ dá
$\int_\Omega\norm{\nabla u}^2 = 0$, de modo que $\nabla u \equiv 0$
e $u$ é constante em cada componente; o \omterm{def:b3:topology:interior}{fecho} de toda componente
encontra $\partial\Omega$ (limitação), onde $u = 0$:
$u \equiv 0$. Duas \omterm{prop:b3:conformal:harmonicholo}{funções harmônicas} com valores de \omterm{def:b3:forms:boundary}{bordo}
iguais diferem por uma tal $u$: elas coincidem --- unicidade
para o problema de Dirichlet, complementando a teoria de existência
no disco do~\cref{ch:b3:conformal}.
\end{solution}

\begin{solution}{exo:b3:forms:9}
Pelo~\cref{thm:b3:forms:pullback}(c),
$\varphi^*(f\,\dd y_1\wedge\dots\wedge\dd y_n) =
(f\circ\varphi)\,\det(D\varphi)\,\dd
x_1\wedge\dots\wedge\dd x_n$, de modo que a identidade de \omterm{def:b3:forms:pullback}{pullback} se lê
\[
\int_U(f\circ\varphi)\,\det(D\varphi)\,\dd x = \int_Vf\,\dd
y .
\]
Como $\det D\varphi > 0$ em toda parte, $\det D\varphi =
\abs{\det D\varphi}$, e isso é literalmente a
fórmula de mudança de variáveis
(\cref{thm:b3:product:changeofvar}) para integrandos \omterm{def:b3:topology:continuity}{contínuos}
de suporte \omterm{def:b3:topology:compact}{compacto}: cada enunciado é o
outro. O valor absoluto foi para a \emph{hipótese}:
a \omterm{def:b3:forms:orientation}{orientação}. Para $\varphi$ que inverte a \omterm{def:b3:forms:orientation}{orientação}, a identidade
entre formas adquire um sinal de menos global (as formas sentem a
\omterm{def:b3:forms:orientation}{orientação}), ao passo que a fórmula de \omterm{def:b3:measure:measure}{medida} mantém $\abs{\det}$
(as \omterm{def:b3:measure:measure}{medidas} não): duas contabilidades de um mesmo jacobiano.
\end{solution}

\begin{solution}{exo:b3:forms:10}
A casca $A = \{r_1 \leq \norm x \leq r_2\}$ é uma
$3$-subvariedade compacta com \omterm{def:b3:forms:boundary}{bordo} $S_{r_2}\cup S_{r_1}$; as
orientações induzidas são a \omterm{def:b3:forms:orientation}{orientação} usual da esfera em
$S_{r_2}$ (para fora de $A$ = para longe de $0$) e a
\emph{oposta} em $S_{r_1}$ (para fora de $A$ = em direção a
$0$). Stokes com $\dd\omega = 0$:
\[
0 = \int_A\dd\omega = \int_{S_{r_2}}\omega -
\int_{S_{r_1}}\omega .
\]
Forma do ângulo sólido: com $\rho = \norm x$ e $\sigma =
x\,\dd y\wedge\dd z + y\,\dd z\wedge\dd x + z\,\dd
x\wedge\dd y$, $\omega = \rho^{-3}\sigma$ e
\[
\dd\omega = \sum_i\partial_i\bigl(x_i\rho^{-3}\bigr)\,
\dd x\wedge\dd y\wedge\dd z
= \bigl(3\rho^{-3} - 3\rho^{-5}\cdot\rho^2\bigr)\dd
x\wedge\dd y\wedge\dd z = 0 .
\]
Em $S_r$, para $v_1, v_2$ tangentes: $\omega(v_1, v_2) =
r^{-3}\det(x, v_1, v_2) = r^{-3}\,r\det(\nu, v_1, v_2) =
r^{-2}\,\dd S(v_1, v_2)$, de modo que $\int_{S_r}\omega =
r^{-2}\cdot4\pi r^2 = 4\pi$: constante em $r$, como a deformação
prevê, e não nula --- de modo que $\omega$ é \omterm{def:b3:forms:closedexact}{fechada}, mas não
\omterm{def:b3:forms:closedexact}{exata}, em $\R^3\setminus\{0\}$ (uma \omterm{def:b3:forms:closedexact}{forma exata} integra
$0$ sobre a $S_r$ sem \omterm{def:b3:forms:boundary}{bordo}, por Stokes). Essa é a lei de Gauss:
o \omterm{ex:b3:ode:planeclassification}{fluxo} do campo de uma carga unitária através de qualquer esfera que a
envolva é $4\pi$, qualquer que seja o raio.
\end{solution}

\begin{solution}{exo:b3:forms:11}
Escreva $\omega = c\,\omega_\theta + \dd f$
(\cref{exo:b3:forms:12}) com $c =
\frac1{2\pi}\int_{S^1}\omega$. Então
\[
\int_\gamma\omega = c\int_\gamma\omega_\theta +
\int_\gamma\dd f
= c\cdot2\pi\operatorname{Ind}_\gamma(0) + 0
= n\int_{S^1}\omega,
\]
pela~\cref{prop:b3:forms:exactloop} e pela definição do
índice. O único inteiro $n$ controla todo período no
plano perfurado: as $1$-formas \omterm{def:b3:forms:closedexact}{fechadas} não conseguem distinguir dois
laços de mesmo número de voltas.
\end{solution}

\begin{solution}{exo:b3:forms:12}
Seja $\alpha = \omega - c\,\omega_\theta$: \omterm{def:b3:forms:closedexact}{fechada}, e
$\int_{S^1}\alpha = 0$ pela escolha de $c$
($\int_{S^1}\omega_\theta = 2\pi$). Puxe para trás pela aplicação
polar $\Phi(\rho, \theta) = (\rho\cos\theta, \rho\sin\theta)$,
um difeomorfismo local sobrejetor $(0, \infty)\times\R \to
U$: $\Phi^*\alpha$ é \omterm{def:b3:forms:closedexact}{fechada}
(\cref{thm:b3:forms:pullback}(b)) no aberto \emph{convexo}
$(0,\infty)\times\R$, logo \omterm{def:b3:forms:closedexact}{exata}
(\cref{thm:b3:forms:poincare}): $\Phi^*\alpha = \dd g$.
Para $\rho$ fixo: $g(\rho, \theta + 2\pi) - g(\rho, \theta)
= \int_\theta^{\theta + 2\pi}\partial_\theta g\,\dd s$ é
a integral de $\alpha$ em torno do círculo de raio
$\rho$, que vale $\int_{S^1}\alpha = 0$ (a coroa
entre os dois círculos é uma superfície compacta com \omterm{def:b3:forms:boundary}{bordo};
Stokes como no~\cref{exo:b3:forms:10}, uma dimensão abaixo).
Logo $g$ é $2\pi$-periódica em $\theta$ e desce a uma
função bem definida $f$ em $U$ com $f\circ\Phi = g$;
$f$ é suave ($\Phi$ é um difeomorfismo local) e
$\Phi^*(\dd f) = \dd g = \Phi^*\alpha$ força $\dd f =
\alpha$. Logo $\omega = c\,\omega_\theta + \dd f$.
Unicidade: integrar sobre $S^1$ fixa $c$, já que as \omterm{def:b3:forms:closedexact}{formas
exatas} têm período nulo; e $f$ é única a menos de constante
aditiva. A aplicação $[\omega] \mapsto
\frac1{2\pi}\int_{S^1}\omega$ é, portanto, um isomorfismo
linear $H^1(\R^2\setminus\{0\}) \to \R$, e a classe
de $\omega_\theta$ gera: o buraco é exatamente
unidimensional, do ponto de vista cohomológico.
\end{solution}

\begin{solution}{pb:b3:forms:1}
\textbf{1.} $\dd\sigma = \sum_i(-1)^{i-1}\,\dd x_i\wedge\dd
x_1\wedge\dots\wedge\widehat{\dd x_i}\wedge\dots\wedge\dd
x_n$; carregar $\dd x_i$ através dos $i-1$ fatores precedentes
custa $(-1)^{i-1}$, o que cancela o prefator: cada um dos
$n$ termos vale $\dd x_1\wedge\dots\wedge\dd x_n$, de modo que
$\dd\sigma = n\,\dd x_1\wedge\dots\wedge\dd x_n$. Stokes na
bola: $\int_S\sigma = \int_{\bar B}\dd\sigma =
n\operatorname{vol}(\bar B) > 0$.

\textbf{2.} $n = 2$: $\sigma = x\,\dd y - y\,\dd x$; em
$\gamma(t) = (\cos t, \sin t)$, $\gamma^*\sigma = (\cos^2t +
\sin^2t)\,\dd t = \dd t$, a forma comprimento de arco:
$\int_S\sigma = 2\pi = 2\operatorname{vol}(\bar B^2)$.
$n = 3$: $\sigma\vert_S$ é a forma de área
(\cref{exo:b3:forms:5} com $\nu(x) = x$):
$\int_S\sigma = 4\pi = 3\cdot\tfrac{4\pi}3$. Ambas batem com
a questão 1.

\textbf{3.} Derivando $\norm\varphi^2 = 1$: $2\langle
D\varphi(x)h, \varphi(x)\rangle = 0$ para todo $h$, de modo que
$\operatorname{im}D\varphi(x) \subseteq \varphi(x)^\perp$, um
hiperplano: $\operatorname{rk}D\varphi(x) \leq n - 1$. Como
$\dd\sigma$ é uma $n$-forma (questão 1), pontualmente
$\varphi^*(\dd\sigma)_x =
(D\varphi(x))^*(\dd\sigma)_{\varphi(x)} = 0$ pela~\cref{prop:b3:forms:linearpullback}(2): uma aplicação com valores na
esfera não tem espaço para puxar de volta uma forma de volume.

\textbf{4.} A falha é o segundo ``logo'': na $S$ de dimensão
$(n-1)$, toda $(n-1)$-forma é trivialmente
\omterm{def:b3:forms:closedexact}{fechada} (não há $n$-formas não nulas numa variedade de dimensão
$(n-1)$), mas \emph{\omterm{def:b3:forms:closedexact}{fechada} não significa \omterm{def:b3:forms:closedexact}{exata}},
e $\int_S\dd\eta = 0$ exige uma primitiva $\eta$ efetivamente
definida em $S$. A restrição de $\sigma$ é precisamente não
\omterm{def:b3:forms:closedexact}{exata} --- sua integral vale $n\operatorname{vol}(\bar B) \neq
0$ --- e essa não exatidão move o problema inteiro.

\textbf{5.} Vamos integrar $r^*\sigma$ sobre a esfera
e contar de duas maneiras. Como $r$ fixa $S$ ponto a ponto, a
integral vale $\int_S\sigma = n\operatorname{vol}(\bar B)
\neq 0$. Como $r$ está definida na bola, Stokes converte
a mesma integral em $\int_{\bar B}\dd(r^*\sigma) =
\int_{\bar B}r^*(\dd\sigma)$; e, como $r$ toma valores
na esfera, a questão 3 anula esse integrando. Um
número, dois valores: a retração não pode existir.

\textbf{6.} Ambas as integrais são calculadas por parametrizações
diretas $\gamma$ de peças de $S$
(\cref{def:b3:forms:integral}); como $r\circ\gamma =
\gamma$ (a parametrização cai em $S$, em que $r$ é a
identidade), $\gamma^*(r^*\sigma) = (r\circ\gamma)^*\sigma =
\gamma^*\sigma$ (\cref{thm:b3:forms:pullback}(a)): os
integrandos locais coincidem, e qualquer \omterm{lem:b3:forms:partition}{partição da unidade} dá
$\int_Sr^*\sigma = \int_S\sigma$.

\textbf{7.} $r^*\sigma$ é uma $(n-1)$-forma suave numa
\omterm{def:b3:topology:topology}{vizinhança} do \omterm{def:b3:topology:compact}{compacto} orientado $\bar B$, cujo
\omterm{def:b3:forms:boundary}{bordo} com a \omterm{def:b3:forms:orientation}{orientação} induzida é $S$: Stokes
(\cref{thm:b3:forms:stokes}) dá exatamente
$\int_Sr^*\sigma = \int_{\bar B}\dd(r^*\sigma)$.

\textbf{8.} $\dd(r^*\sigma) = r^*(\dd\sigma)$
(\cref{thm:b3:forms:pullback}(b)), que se anula pela
questão 3 aplicada a $\varphi = r$. Encadeando as questões
6--8:
\[
0 < n\operatorname{vol}(\bar B) = \int_S\sigma =
\int_Sr^*\sigma = \int_{\bar B}\dd(r^*\sigma) = 0 :
\]
absurdo. \emph{Não existe retração suave de $\bar B^n$
sobre $S^{n-1}$} ($n \geq 2$).

\textbf{9.} Uma $r\colon\intcc{-1}1\to\{-1,1\}$ \omterm{def:b3:topology:continuity}{contínua}
com $r(\pm1) = \pm1$ levaria um \omterm{def:b3:topology:connected}{conexo} sobre o
desconexo $\{-1, 1\}$, impossível: imagens \omterm{def:b3:topology:continuity}{contínuas} de
\omterm{def:b3:topology:connected}{conexos} são \omterm{def:b3:topology:connected}{conexas} --- equivalentemente, o
teorema do valor intermediário forçaria $r$ a assumir o valor
$0$. O mesmo enunciado em toda dimensão é exatamente a Parte
II; a \omterm{def:b3:topology:connected}{conexidade} é a sombra $1$-dimensional do
obstáculo cohomológico $\int_S\sigma \neq 0$.

\textbf{10.} $x \mapsto \norm{g(x) - x}$ é \omterm{def:b3:topology:continuity}{contínua} e
$> 0$ em toda parte no \omterm{def:b3:topology:compact}{compacto} $\bar B$: seu mínimo
$\delta$ é atingido, logo $> 0$.

\textbf{11.} $\norm{x + tu}^2 = 1$ se lê $t^2 + 2t\langle x,
u\rangle + \norm x^2 - 1 = 0$, cujas raízes são
\[
t_\pm = -\langle x, u\rangle \pm \sqrt{\langle x,
u\rangle^2 + 1 - \norm x^2} .
\]
Seu produto é $\norm x^2 - 1 \leq 0$: as raízes ficam de lados opostos de
$0$ (ou uma delas se anula), de modo que o parâmetro do raio --- a
raiz não negativa --- é $t(x) = t_+$.

\textbf{12.} Se o radicando se anulasse em $x \in \bar B$:
sendo ambos os termos não negativos, $\norm x = 1$ e
$\langle x, u(x)\rangle = 0$, isto é, $\langle x, x -
g(x)\rangle = 0$, isto é, $\langle x, g(x)\rangle = 1$. Por
Cauchy--Schwarz, $1 = \langle x, g(x)\rangle \leq
\norm x\,\norm{g(x)} \leq 1$: igualdade em toda a cadeia, o que
força $g(x)$ colinear com $x$, de norma $1$, positivamente:
$g(x) = x$ --- excluído. Logo o radicando é \omterm{def:b3:topology:continuity}{contínuo} e
$> 0$ em $\bar B$, logo minorado ali por algum $c > 0$
e numa \omterm{def:b3:topology:topology}{vizinhança} (\omterm{def:b3:topology:continuity}{continuidade} uniforme). Nessa
\omterm{def:b3:topology:topology}{vizinhança}, $u$ é suave ($\norm{x - g(x)} \geq
\delta/2$, encolhendo se necessário), o radicando permanece $\geq
c/2$, e a raiz quadrada é suave em $\intoo0\infty$: $r$
é suave perto de $\bar B$.

\textbf{13.} $\norm{r(x)} = 1$ por construção de $t(x)$:
$r(\bar B) \subseteq S$. Para $x \in S$: $\langle x,
u(x)\rangle = \frac{1 - \langle x, g(x)\rangle}{\norm{x -
g(x)}} \geq 0$ (Cauchy--Schwarz mais uma vez), e $\norm x =
1$ reduz o radicando a $\langle x, u\rangle^2$, cuja
raiz quadrada é o próprio $\langle x, u\rangle$ (ele é $\geq
0$): $t(x) = 0$ e $r(x) = x$. Logo $r$ é uma retração suave
da bola sobre a esfera --- contradizendo
a Parte II. \emph{Toda autoaplicação suave de $\bar B$ tem
ponto fixo.}

\textbf{14.} $\varepsilon > 0$ exatamente como na questão 10.
Os polinômios formam uma subálgebra de $\mathcal C(\bar B,
\R)$ que contém as constantes e separa pontos (as $x
\mapsto x_i$ separam), de modo que Stone--Weierstrass
(\cref{thm:b3:complete:stoneweierstrass}) aproxima cada
coordenada: tome polinômios $p_i$ com $\sup_{\bar
B}\abs{p_i - f_i} < \frac{\varepsilon}{2\sqrt n}$; a
aplicação vetorial $p = (p_1, \dots, p_n)$ satisfaz então
$\sup_{\bar B}\norm{p - f} \leq
\bigl(\sum_i\sup_{\bar B}\abs{p_i - f_i}^2\bigr)^{1/2} <
\varepsilon/2$.

\textbf{15.} Em $\bar B$: $\norm p \leq \norm f +
\frac\varepsilon2 \leq 1 + \frac\varepsilon2$, de modo que
$\norm{g} = \frac{\norm p}{1 + \varepsilon/2} \leq 1$:
$g(\bar B) \subseteq \bar B$. Além disso, $\norm{g - p} =
\frac{\varepsilon/2}{1 + \varepsilon/2}\norm p \leq
\frac\varepsilon2$, logo $\norm{g - f} \leq \norm{g - p} +
\norm{p - f} < \varepsilon$ em $\bar B$.

\textbf{16.} $g$ é polinomial, logo suave, e leva
$\bar B$ em si mesmo: a Parte III fornece $x_0 = g(x_0)$.
Então $\norm{f(x_0) - x_0} = \norm{f(x_0) - g(x_0)} <
\varepsilon = \min_{\bar B}\norm{f - \operatorname{id}}$:
contradição. \emph{Toda \omterm{def:b3:topology:continuity}{aplicação contínua} $\bar B^n \to
\bar B^n$ tem ponto fixo.}

\textbf{17.} \emph{Bola aberta:} $f(x) = \frac{x + e_1}2$ leva
$B$ em $B$ ($\norm{f(x)} < 1$ estritamente) e seu único ponto
fixo $e_1$ está na esfera: perdeu-se a \omterm{def:b3:topology:compact}{compacidade}.
\emph{Esfera:} a aplicação antipodal $x \mapsto -x$ não tem
ponto fixo; $S$ é compacta, mas tem a \omterm{def:b3:topology:topology}{topologia} ``errada''
--- ela é exatamente a não retração da Parte II.
\emph{Coroa:} uma rotação de qualquer ângulo $\not\equiv 0$ nada
fixa; o buraco abriga a rotação --- perdeu-se a convexidade (mais
precisamente, a \omterm{def:b3:topology:topology}{topologia} do tipo bola). O teorema de Brouwer
é, na verdade, sobre conjuntos \emph{\omterm{def:b3:topology:compact}{compactos} convexos}, como a questão 18
explora.

\textbf{18.} Para $x \in \Delta$: algum $x_j > 0$, de modo que $(Ax)_i
\geq A_{ij}x_j > 0$ para todo $i$; logo $\norm{Ax}_1 > 0$
e $T(x) = Ax/\norm{Ax}_1$ está bem definida, é \omterm{def:b3:topology:continuity}{contínua} e
cai em $\Delta$ (entradas positivas de soma $1$).
$\Delta$ é convexo, \omterm{def:b3:topology:compact}{compacto}, com \omterm{def:b3:topology:interior}{interior} não vazio no
hiperplano afim $\{\sum x_i = 1\} \cong \R^{n-1}$; a
aplicação radial a partir de seu baricentro --- cada raio a partir do
baricentro encontra $\partial\Delta$ em exatamente um ponto, por
convexidade e \omterm{def:b3:topology:compact}{compacidade}, e a função de calibre correspondente
é \omterm{def:b3:topology:continuity}{contínua} --- é um \omterm{def:b3:topology:continuity}{homeomorfismo} $\Delta \to
\bar B^{n-1}$. Transportando Brouwer por ele: $T$ tem
ponto fixo $x^*$, isto é, $Ax^* = \lambda x^*$ com $\lambda
= \norm{Ax^*}_1 > 0$; e $x^* = Ax^*/\lambda$ tem entradas estritamente
positivas pelo cálculo inicial. \emph{Uma
matriz positiva tem autovetor positivo.}

\textbf{19.} A questão 18 dá $A\pi = \lambda\pi$, $\pi \in
\Delta$, $\pi > 0$. Some as coordenadas: $\sum_i(A\pi)_i =
\sum_j\pi_j\sum_iA_{ij} = \sum_j\pi_j = 1$ (as colunas somam
$1$), ao passo que $\sum_i\lambda\pi_i = \lambda$. Logo $\lambda = 1$
e $A\pi = \pi$: um vetor de probabilidade estacionário --- o
equilíbrio da cadeia de Markov, o núcleo matemático do
PageRank.

\textbf{20.} Em $S$: $\norm{F_t(x)}^2 = \norm x^2 +
2t\langle x, v(x)\rangle + t^2\norm{v(x)}^2 = 1 + 0 + t^2$
pela tangência e $\norm v = 1$: $F_t(S) \subseteq
\sqrt{1+t^2}\,S$.

\textbf{21.} Fixe um atlas finito de parametrizações diretas
$\gamma$ e uma \omterm{lem:b3:forms:partition}{partição da unidade}, independentes de $t$. Numa
carta, $F_t\circ\gamma = \gamma + t(v\circ\gamma)$, de modo que cada
coeficiente de $(F_t\circ\gamma)^*\sigma$ é uma soma de
produtos de um fator $(x_i + tv_i)\circ\gamma$ (afim em
$t$) por um determinante $(n-1)\times(n-1)$ com entradas
afins em $t$: um polinômio em $t$ de grau $\leq n$ com
coeficientes suaves na variável da carta. Multiplicando pelas
funções da partição, independentes de $t$, e integrando
termo a termo: $P(t) = \sum_{k=0}^n c_kt^k$, um polinômio.

\textbf{22.} \emph{Injetividade:} $v$ é lipschitziana em $S$
(suave num \omterm{def:b3:topology:compact}{compacto}), digamos de constante $L$; para $\abs t <
1/L$, $\norm{F_t(x) - F_t(y)} \geq (1 - \abs
tL)\norm{x - y} > 0$. \emph{Difeomorfismo:} $G_t =
F_t/\sqrt{1 + t^2}$ leva $S$ em $S$; em cartas, seus jacobianos
convergem uniformemente aos de $G_0 = \operatorname{id}$ quando
$t \to 0$, de modo que, para $t$ pequeno, eles são inversíveis e $G_t$ é
um difeomorfismo local (\cref{thm:b3:submanifolds:ift} em
cartas), injetor, de imagem aberta (difeomorfismo local) e
compacta na \emph{\omterm{def:b3:topology:connected}{conexa}} $S^{n-1}$ ($n \geq 2$):
imagem $= S$, de modo que $G_t$ é um difeomorfismo de $S$.
\emph{Escala:} sob $s_\lambda(x) = \lambda x$, cada
coeficiente $x_i$ ganha $\lambda$ e cada uma das $n - 1$
diferenciais ganha $\lambda$: $s_\lambda^*\sigma =
\lambda^n\sigma$. Como $F_t = s_{\sqrt{1+t^2}}\circ G_t$:
\[
P(t) = \int_SG_t^*\bigl(s_{\sqrt{1+t^2}}^{\,*}\sigma\bigr)
= (1 + t^2)^{n/2}\int_SG_t^*\sigma
= (1 + t^2)^{n/2}\int_S\sigma
\quad\text{para } t,
\]
a última igualdade pelo~\cref{lem:b3:forms:consistency} ($G_t$
é um difeomorfismo de $S$, que preserva a \omterm{def:b3:forms:orientation}{orientação} para $t$
pequeno: os determinantes jacobianos em cartas variam \omterm{def:b3:topology:continuity}{continuamente},
nunca se anulam e são positivos em $t = 0$).

\textbf{23.} Se $n$ é ímpar e existe um campo tangente unitário
suave, as questões 21--22 fazem o polinômio
$P(t)/\int_S\sigma$ (legítimo: $\int_S\sigma \neq 0$ pela
questão 1) coincidir com $(1 + t^2)^{n/2}$ perto de $0$; duas
funções suaves que coincidem perto de $0$, sendo uma delas polinomial,
forçam $(1 + t^2)^{n/2}$ a ser esse polinômio em todo
$\R$. Mas, se $Q(t)^2 = (1 + t^2)^n$ com $Q \in \R[t]$,
a fatoração única em $\R[t]$
(\cref{ch:b3:rings}) dá ao \omterm{def:b3:rings:divisibility}{irredutível} $t^2 + 1$ multiplicidade
par em $Q^2$ e multiplicidade ímpar $n$ em
$(1 + t^2)^n$: impossível. Logo não existe campo tangente unitário
suave em $S^{n-1}$ para $n$ ímpar --- as esferas de
dimensão par. Por fim, um campo tangente meramente \omterm{def:b3:topology:continuity}{contínuo} e
nunca nulo $w$ produziria um: estenda-o a uma
\omterm{def:b3:topology:topology}{vizinhança} por $\tilde w(x) = w(x/\norm x)$, regularize
componente a componente (\cref{thm:b3:lp:regularization}) até um
$v_0$ suave com $\sup_S\norm{v_0 - \tilde w} <
\frac12\min_S\norm w$, projete tangencialmente $v_1(x) = v_0(x)
- \langle v_0(x), x\rangle x$ --- em $S$ isso altera $v_0$
no máximo por sua componente normal, ela própria no máximo $\norm{v_0 -
\tilde w}$, já que $\tilde w$ é tangente, de modo que $\norm{v_1 - \tilde
w} \leq 2\norm{v_0 - \tilde w} < \min\norm w$ e $v_1$
nunca se anula em $S$ --- e normalize: $v =
v_1/\norm{v_1}$ é um campo tangente unitário suave. Logo, em
toda esfera de dimensão par, \emph{todo campo tangente \omterm{def:b3:topology:continuity}{contínuo}
tem um zero}: todo vento na Terra deixa um ponto de
calmaria.

\textbf{24.} $\langle v(x), x\rangle = \sum_j(-y_jx_j +
x_jy_j) = 0$: tangente; $\norm{v(x)} = \norm x = 1$: unitário.
A suavidade é clara (aplicação linear). Logo toda esfera de dimensão ímpar
carrega um campo tangente unitário suave --- a multiplicação
por $\iu$ ao longo das retas complexas --- e o obstáculo da questão 23
é exatamente a paridade da dimensão. No
argumento polinomial, para $n$ par a função $(1 +
t^2)^{n/2}$ \emph{é} um polinômio, e nenhuma contradição
surge: a demonstração não apenas deixa de se aplicar --- sua
conclusão é genuinamente falsa, como $v$ testemunha.

\textbf{25.} Suponha que $f$ nunca se anule em $\bar B$. Então
$g(x) = -f(x)/\norm{f(x)}$ é \omterm{def:b3:topology:continuity}{contínua} $\bar B \to S
\subseteq \bar B$, e Brouwer (questão 16) fornece $x^* =
g(x^*)$. Como $g$ toma valores em $S$, $x^* \in S$, e
\[
\langle f(x^*), x^*\rangle = \bigl\langle f(x^*),\
-\tfrac{f(x^*)}{\norm{f(x^*)}}\bigr\rangle
= -\norm{f(x^*)} < 0,
\]
contradizendo a hipótese de \omterm{def:b3:forms:boundary}{bordo}. Logo $f$ tem um zero.
Sobrejetividade: dado $y \in \R^n$, aplique o que precede a $f(x) =
F(x) - y$ numa bola $\bar B(0, R)$ com $R$ tão grande que
$\langle F(x), x\rangle \geq \norm y\,\norm x$ na esfera
de raio $R$ (coercividade); então $\langle f(x), x\rangle =
\langle F(x), x\rangle - \langle y, x\rangle \geq 0$ ali
(Cauchy--Schwarz), e o enunciado reescalado dá um zero
de $f$: $F(x) = y$. Todo campo \omterm{def:b3:topology:continuity}{contínuo} coercivo é sobrejetor
--- a sombra sem grau dos teoremas variacionais de
existência, entregue pela \omterm{def:b3:topology:topology}{topologia} pura.
\end{solution}
